\documentclass{article}
\usepackage[margin=1in]{geometry}
\usepackage{mathtools, amsfonts, amsthm, xcolor, listings}

\newcommand{\R}{\mathbb{R}}
\newcommand{\N}{\mathbb{N}}
\newcommand{\I}{\mathbb{I}}
\newcommand{\Q}{\mathbb{Q}}
\newcommand{\Z}{\mathbb{Z}}

\newcommand{\set}[1]{\left\{#1\right\}}
\newcommand{\ang}[1]{\left\langle #1 \right\rangle}
\newcommand{\inv}{^{-1}}

\newcommand{\normsub}{\lhd}

\newcommand{\reflection}[5]{
    \begin{itemize}
    \item Points attempted: #1
    \item Self-assessed score: #2/#1.
    \item Time spent: #3 minutes.
    \item Difficulty: #4/10.
    \item Questions/Feedback: #5.
    \end{itemize}
}


\title{HW 12}
\author{Sam Ly}

\begin{document}
\maketitle

https://github.com/samlyme/Assignments/tree/cs4080/ch24/main/CS4080

\section*{Chapter 24}

\begin{enumerate}
    \item {
        Reading and writing the \texttt{ip} field is one of the most frequent operations
        inside the bytecode loop. Implement the optimization of caching it in a local
        C variable. Write a couple of benchmarks and say whether the extra complexity
        is worth it.

        The main change is to keep the current instruction pointer in a local variable
        inside \texttt{run()}:

        \begin{verbatim}
static InterpretResult run() {
  CallFrame* frame = &vm.frames[vm.frameCount - 1];
  register uint8_t* ip = frame->ip;

#define READ_BYTE() (*ip++)
#define READ_SHORT() \
    (ip += 2, (uint16_t)((ip[-2] << 8) | ip[-1]))
        \end{verbatim}

        Any opcode that changes control flow now changes the local \texttt{ip}, not
        \texttt{frame->ip}. For example, \texttt{OP\_JUMP} adds to \texttt{ip} and
        \texttt{OP\_LOOP} subtracts from it. The important bookkeeping is around
        calls, returns, and errors. Before calling another function, the VM stores the
        local value back into the current frame, then reloads \texttt{ip} from the new
        top frame after the call:

        \begin{verbatim}
case OP_CALL: {
  int argCount = READ_BYTE();
  frame->ip = ip;
  if (!callValue(peek(argCount), argCount)) {
    return INTERPRET_RUNTIME_ERROR;
  }
  frame = &vm.frames[vm.frameCount - 1];
  ip = frame->ip;
  break;
}
        \end{verbatim}

        The same reload is needed after \texttt{OP\_RETURN}. Runtime error paths also
        need \texttt{frame->ip = ip;} immediately before calling \texttt{runtimeError()},
        otherwise the line number reported for the error can be wrong.

        I would benchmark one call-heavy program like recursive Fibonacci and one
        loop-heavy program that mostly executes arithmetic and jumps. This optimization
        should help both, but it especially helps the call-heavy case because
        \texttt{ip} is touched constantly while dispatching bytecode. A single digit
        speedup is enough to be worth it here because the added complexity is localized
        to the VM dispatch loop and does not affect the language semantics.
    }

    \item {
        Native function calls are fast because they do not validate the number of
        arguments. Add arity checking.

        I would make native functions carry their expected arity the same way
        \texttt{ObjFunction} already does. Instead of storing only the function pointer
        in the native object, store both the C function and an integer arity:

        \begin{verbatim}
typedef struct {
  Obj obj;
  NativeFn function;
  int arity;
} ObjNative;
        \end{verbatim}

        Then \texttt{defineNative()} takes an arity argument, and \texttt{callValue()}
        checks it before invoking the native:

        \begin{verbatim}
case OBJ_NATIVE: {
  ObjNative* native = AS_NATIVE(callee);
  if (argCount != native->arity) {
    runtimeError("Expected %d arguments but got %d.",
                 native->arity, argCount);
    return false;
  }

  Value result = native->function(argCount, vm.stackTop - argCount);
  vm.stackTop -= argCount + 1;
  push(result);
  return true;
}
        \end{verbatim}

        This adds one integer comparison to every native call. That is a reasonable
        cost because it prevents native code from reading missing arguments out of
        uninitialized stack slots.
    }

    \item {
        Extend the native function system so a native function can signal a runtime
        error. How does this affect native call performance?

        The cleanest design is to separate the success signal from the returned Lox
        value. I would change \texttt{NativeFn} to return a boolean and have the native
        write either its result or an error message into the stack slot before the
        first argument:

        \begin{verbatim}
typedef bool (*NativeFn)(int argCount, Value* args);

static bool clockNative(int argCount, Value* args) {
  args[-1] = NUMBER_VAL((double)clock() / CLOCKS_PER_SEC);
  return true;
}

static bool sqrtNative(int argCount, Value* args) {
  if (!IS_NUMBER(args[0])) {
    args[-1] = OBJ_VAL(copyString("sqrt() expects a number.", 24));
    return false;
  }

  args[-1] = NUMBER_VAL(sqrt(AS_NUMBER(args[0])));
  return true;
}
        \end{verbatim}

        The VM then keeps the return value in the callee's old stack slot on success,
        or prints the message from that slot on failure:

        \begin{verbatim}
if (native->function(argCount, vm.stackTop - argCount)) {
  vm.stackTop -= argCount;
  return true;
}

runtimeError(AS_CSTRING(vm.stackTop[-argCount - 1]));
return false;
        \end{verbatim}

        This adds a branch to every native call, so a native-call microbenchmark would
        get slower. I still think it is the right tradeoff because a real standard
        library needs to reject bad dynamic values safely instead of letting C code
        crash or read invalid memory.
    }

    \item {
        Add more useful native functions. What did you add, and how do they affect the
        feel of the language?

        I would add a small math and utility base:

        \begin{itemize}
            \item \texttt{sqrt(n)} for square roots.
            \item \texttt{floor(n)} for rounding down.
            \item \texttt{pow(a, b)} for exponentiation.
            \item \texttt{str(value)} for converting simple values to strings.
            \item \texttt{type(value)} for debugging dynamic values.
        \end{itemize}

        Example Lox program:

        \begin{verbatim}
fun distance(x1, y1, x2, y2) {
  var dx = x2 - x1;
  var dy = y2 - y1;
  return sqrt(pow(dx, 2) + pow(dy, 2));
}

print distance(0, 0, 3, 4);
print "answer is " + str(42);
print type(clock());
        \end{verbatim}

        These do not make the language more expressive in theory, but they make it much
        more practical. Without them, even basic numerical programs require awkward
        hand-written helpers. The main design pressure is that each native function
        needs clear arity and runtime type errors so the boundary between dynamic Lox
        and static C stays safe.
    }
\end{enumerate}

\section*{Chapter 25}
https://github.com/samlyme/Assignments/tree/cs4080/ch25/main/CS4080
\begin{enumerate}
    \item {
        Change clox so it only wraps functions in \texttt{ObjClosure}s when the function
        actually needs upvalues. Compare the complexity and performance with always
        wrapping functions.

        The compiler can choose between \texttt{OP\_CLOSURE} and \texttt{OP\_CONSTANT}
        after \texttt{endCompiler()} returns the completed \texttt{ObjFunction}:

        \begin{verbatim}
ObjFunction* function = endCompiler();
uint8_t constant = makeConstant(OBJ_VAL(function));

if (function->upvalueCount > 0) {
  emitBytes(OP_CLOSURE, constant);
  for (int i = 0; i < function->upvalueCount; i++) {
    emitByte(compiler.upvalues[i].isLocal ? 1 : 0);
    emitByte(compiler.upvalues[i].index);
  }
} else {
  emitBytes(OP_CONSTANT, constant);
}
        \end{verbatim}

        The VM then has to accept either \texttt{ObjFunction} or \texttt{ObjClosure} in a
        call frame. That means the frame can no longer store only a closure pointer.
        It either stores a generic \texttt{Obj*}, or it stores both the callable object
        and the underlying \texttt{ObjFunction}. The latter keeps bytecode access simple:

        \begin{verbatim}
typedef struct {
  Obj* callable;
  ObjFunction* function;
  uint8_t* ip;
  Value* slots;
} CallFrame;
        \end{verbatim}

        The performance tradeoff is mixed. Programs with many plain functions avoid
        allocating unnecessary closure objects, so they can become faster and allocate
        less. Programs that heavily use closures do not gain much, and the VM grows
        more complicated because call sites and stack traces have to handle two
        callable representations. I would weight ordinary non-closure function calls
        more heavily than a synthetic closure-allocation benchmark, but I still would
        want real program benchmarks before committing to the optimization. My default
        choice would be to keep always wrapping functions unless the allocation cost
        shows up clearly in profiling.
    }

    \item {
        How should Lox behave when closures capture a loop variable? Change the
        implementation to create a new variable for each loop iteration.

        I think each iteration should get a fresh variable. This matches what
        programmers usually mean when they create closures inside a loop: each closure
        should remember the value from the iteration that created it.

        The compiler can implement this by introducing a hidden inner scope around the
        loop body when the \texttt{for} loop declares a variable:

        \begin{enumerate}
            \item Remember the original loop variable's name and stack slot.
            \item Before compiling the body, begin a new scope and create a local with
            the same name initialized from the outer loop variable.
            \item Compile the body so closures resolve to the inner per-iteration
            variable.
            \item After the body, copy the inner variable back to the outer loop
            variable so assignments in the body still affect the loop condition and
            increment.
            \item End the inner scope, closing the captured variable if a closure
            refers to it.
        \end{enumerate}

        Conceptually, this makes code like this:

        \begin{verbatim}
for (var i = 1; i <= 2; i = i + 1) {
  fun closure() {
    print i;
  }
  save(closure);
}
        \end{verbatim}

        behave more like each loop body starts with a fresh shadowing variable:

        \begin{verbatim}
{
  var i = outerI;
  fun closure() {
    print i;
  }
  save(closure);
  outerI = i;
}
        \end{verbatim}

        With this behavior, two closures created by two different iterations print
        \texttt{1} and \texttt{2}, not the same final value.
    }

    \item {
        Using closures, write a Lox program that models two-dimensional vector
        objects. It should have a constructor, coordinate accessors, and an addition
        method.

        A simple way to model objects with closures is to return a dispatcher function.
        The dispatcher's closed-over variables are the object's fields, and messages
        select which method to return.

        \begin{verbatim}
fun vector(x, y) {
  fun object(message) {
    fun add(other) {
      return vector(x + other("x"), y + other("y"));
    }

    if (message == "x") return x;
    if (message == "y") return y;
    if (message == "add") return add;

    print "unknown message";
    return nil;
  }

  return object;
}

var a = vector(1, 2);
var b = vector(3, 4);
var c = a("add")(b);

print c("x");
print c("y");
        \end{verbatim}

        The constructor \texttt{vector(x, y)} returns a closure that captures
        \texttt{x} and \texttt{y}. Calling the closure with \texttt{"x"} or \texttt{"y"}
        reads a coordinate. Calling it with \texttt{"add"} returns another closure
        that takes the other vector and constructs the sum.
    }
\end{enumerate}

\end{document}
